\documentclass[11pt]{article}

% ============================================================
% GA-SIM-01 — Simulación: priorización y negociación entre roles
% Ingeniería de Requerimientos (ISR-401) · UTEQ · SPA 2026-2027
% Compilación: pdflatex -> bibtex -> pdflatex -> pdflatex
% ============================================================

\usepackage[utf8]{inputenc}
\usepackage[T1]{fontenc}
% babel no disponible en este entorno
\usepackage[scaled]{helvet}
\renewcommand{\familydefault}{\sfdefault}
\usepackage[a4paper,top=2.6cm,bottom=2.3cm,left=2cm,right=2cm,headheight=34pt]{geometry}
\usepackage{graphicx}
\usepackage{array}
\usepackage{tabularx}
\usepackage{longtable}
\usepackage{pdflscape}
\usepackage{colortbl}
\usepackage{xcolor}
\usepackage{enumitem}
\usepackage{fancyhdr}
\usepackage{lastpage}
\usepackage[most]{tcolorbox}
\usepackage{amsmath}
\usepackage{amssymb}
\usepackage[numbers,sort&compress]{natbib}
\usepackage{url}
\usepackage{needspace}
\usepackage{caption}
\usepackage[hidelinks]{hyperref}

% ---------- Textos en español (sin babel) ----------
\renewcommand{\tablename}{Tabla}
\renewcommand{\refname}{Referencias}

% ---------- Paleta institucional ----------
\definecolor{utegreen}{HTML}{0B7A3B}
\definecolor{utegreendk}{HTML}{075C2C}
\definecolor{midblue}{HTML}{1F5C8B}
\definecolor{gold}{HTML}{D4A017}
\definecolor{rowgray}{HTML}{EFEFEF}
\definecolor{boxgray}{HTML}{F4F6F4}
\definecolor{lightred}{HTML}{FBEAEA}
\definecolor{lightgold}{HTML}{FBF3DC}
\definecolor{lightblue}{HTML}{EAF1F7}
\definecolor{lightgreen}{HTML}{DCEFE3}

% ---------- Columnas personalizadas ----------
\newcolumntype{L}[1]{>{\raggedright\arraybackslash}p{#1}}
\newcolumntype{C}[1]{>{\centering\arraybackslash}p{#1}}
\newcolumntype{Y}{>{\raggedright\arraybackslash}X}

% ---------- Encabezado y pie ----------
\pagestyle{fancy}
\fancyhf{}
\lhead{\footnotesize UTEQ \,$\cdot$\, Ingeniería de Requerimientos (ISR-401)}
\rhead{\footnotesize GA-SIM-01 \,$\cdot$\, SPA 2026--2027}
\lfoot{\footnotesize Prof. Gleiston C. Guerrero-Ulloa, PhD}
\rfoot{\footnotesize Página \thepage\ de \pageref{LastPage}}
\renewcommand{\headrulewidth}{0.6pt}
\renewcommand{\footrulewidth}{0.4pt}

\setlength{\parindent}{0pt}
\setlength{\parskip}{5pt}
\setlength{\emergencystretch}{3em}
\renewcommand{\arraystretch}{1.25}

% ---------- Atajos ----------
\newcommand{\thdr}[1]{\cellcolor{utegreen}\textcolor{white}{\textbf{#1}}}
\newcommand{\seccion}[1]{\section{\textcolor{utegreendk}{#1}}}
\newcommand{\subseccion}[1]{\subsection{\textcolor{midblue}{#1}}}
\renewcommand{\thesection}{\arabic{section}.}
\renewcommand{\thesubsection}{\arabic{section}.\arabic{subsection}.}

\newtcolorbox{cajaclave}[1][]{colback=boxgray,colframe=utegreen,boxrule=0.8pt,arc=2mm,fonttitle=\bfseries,breakable,#1}
\newtcolorbox{cajapiso}[1][]{colback=lightred,colframe=red!60!black,boxrule=1pt,arc=2mm,fonttitle=\bfseries,breakable,#1}
\newtcolorbox{cajaaviso}[1][]{colback=lightgold,colframe=gold!80!black,boxrule=0.8pt,arc=2mm,fonttitle=\bfseries,breakable,#1}

\begin{document}

\noindent\includegraphics[width=\linewidth]{logoband.png}

\vspace{4pt}
\begin{center}
	{\small Facultad de Ciencias de la Computación \,\textcolor{gold}{$\bullet$}\, Carrera de Software \,\textcolor{gold}{$\bullet$}\, Ingeniería de Requerimientos (ISR-401)}\\[8pt]
	{\LARGE\textbf{\textcolor{utegreendk}{GUÍA Y RÚBRICA DE EVALUACIÓN}}}\\[4pt]
	{\large\textbf{\textcolor{midblue}{ACTIVIDAD FORMATIVA}}}\\[6pt]
	{\Large\textbf{GA-SIM-01}}\\[2pt]
	{\large Simulación: priorización y negociación entre roles}
\end{center}

\vspace{2pt}
\noindent\textcolor{gold}{\rule{\linewidth}{1.5pt}}

\seccion{Datos generales}

\noindent\begin{tabularx}{\linewidth}{|L{0.30\linewidth}|Y|}
	\hline
	\thdr{Campo} & \thdr{Detalle}\\
	\hline
	Asignatura & Ingeniería de Requerimientos (ISR-401) \,$\cdot$\, 4.\textsuperscript{o} nivel \,$\cdot$\, Paralelo A\\
	\hline
	\rowcolor{rowgray} Docente & Gleiston Cicerón Guerrero Ulloa, PhD \,$\cdot$\, \texttt{gguerrero@uteq.edu.ec}\\
	\hline
	Unidad & Unidad 2. Elicitación y Análisis de Requisitos\\
	\hline
	\rowcolor{rowgray} Resultado de aprendizaje & Aplica técnicas de elicitación y análisis de requisitos para identificar, clasificar y modelar requisitos.\\
	\hline
	Carácter & Formativo: sin calificación; su entrega válida es requisito de la sumativa vinculada\\
	\hline
	\rowcolor{rowgray} Prepara (sumativas vinculadas) & GA-PE-02, EV-PRAC-02\\
	\hline
	Modalidad & Grupal (equipo del PFC)\\
	\hline
	\rowcolor{rowgray} Semana & 5 (23 de noviembre -- 29 de noviembre de 2026)\\
	\hline
	Publicación de la guía & lunes 23 de noviembre de 2026\\
	\hline
	\rowcolor{rowgray} Cierre & sábado 28 de noviembre de 2026, a las 23:59\\
	\hline
	Repositorio & Repositorio del equipo \texttt{ISR401-<SIGLA>}\\
	\hline
\end{tabularx}

\seccion{Propósito}

Experimentar en clase un conflicto de prioridades entre cliente, analista y desarrollador, y resolverlo con técnicas de priorización y negociación ganar-ganar.

\seccion{Fundamento}

Los conflictos de requisitos pueden ser de datos, de intereses, de valores, de relación o estructurales, y se resuelven con estrategias como el acuerdo, el compromiso, la votación o la variante de solución \citep{pohl_2025}. MoSCoW y la votación ponderada hacen explícitas las prioridades \citep{wiegers_beatty_2013}.

\seccion{Desarrollo de la actividad}

\begin{enumerate}[leftmargin=20pt,itemsep=4pt]
	\item \textbf{Roles.} Dentro de cada equipo del PFC se forma un grupo de tres (si el equipo tiene 4 o 5 integrantes, los demás actúan como observadores y ayudan al analista a registrar los conflictos). Cada integrante del grupo recibe del docente, en clase, una tarjeta impresa con su rol (cliente, analista, desarrollador) y sus objetivos secretos. El caso, los 12 requisitos y las plantillas de registro están en el Anexo A.
	\item \textbf{Priorización individual.} Cada rol prioriza 12 requisitos con MoSCoW y reparte 100 puntos.
	\item \textbf{Negociación.} En 20 minutos, el grupo negocia una lista única; el analista registra los conflictos y cómo se resolvieron.
	\item \textbf{Reflexión.} Cada integrante escribe qué tipo de conflicto apareció y qué estrategia funcionó.
\end{enumerate}

\seccion{Productos entregables}

{\small
\begin{longtable}{|L{0.27\linewidth}|L{0.43\linewidth}|L{0.21\linewidth}|}
	\caption{Productos que deben estar en el repositorio al cierre.}\\
	\hline
	\thdr{Ruta en el repositorio} & \thdr{Contenido} & \thdr{Mínimo exigido}\\
	\hline
	\endfirsthead
	\hline
	\thdr{Ruta en el repositorio} & \thdr{Contenido} & \thdr{Mínimo exigido}\\
	\hline
	\endhead
	\texttt{GA-SIM-01/\allowbreak{}main.tex y main.pdf} & Priorizaciones individuales, lista acordada, registro de conflictos y reflexiones & 1\\
	\hline
\end{longtable}
}

\seccion{Normas de entrega}

\begin{enumerate}[leftmargin=20pt,itemsep=3pt]
	\item \textbf{Repositorio.} Todo el trabajo se desarrolla en el repositorio del equipo \texttt{ISR401-<SIGLA>}, que debe ser público en GitHub. Allí deben estar los fuentes (\texttt{.tex}, \texttt{.bib}, figuras, fuentes editables de los modelos y evidencias) y el PDF compilado.
	\item \textbf{Documento en el SGA.} En la tarea del SGA se sube solamente un PDF de \textbf{una sola página} con los datos informativos: estudiante o estudiantes, código y nombre de la tarea, fecha límite de entrega y la URL del repositorio escrita completa en una sola línea. No se sube ningún otro archivo.
	\item \textbf{\LaTeX{} reproducible.} El PDF debe poder regenerarse clonando el repositorio y compilando el archivo \texttt{.tex}. El \texttt{README.md} documenta compilador, archivo principal, orden de comandos y dependencias.
	\item \textbf{Referencias.} Estilo IEEE. Toda referencia debe existir y sus datos (autores, año, título, revista o editorial, volumen, número, páginas y DOI o ISBN) deben corresponder exactamente a la obra real.
	\item \textbf{Plazo.} El cierre es el sábado 28 de noviembre de 2026, a las 23:59. Se evalúa el contenido del repositorio hasta ese momento; lo agregado después no se considera.
	\item \textbf{Trabajo anticipado.} El trabajo se realiza durante la semana, no a última hora; el historial de commits debe mostrarlo. La lista de cotejo registra la constancia observada.
	\item \textbf{Datos personales.} Ningún dato personal de terceros (cédulas, teléfonos, correos personales, audios o videos de participantes, consentimientos firmados) se publica en el repositorio. Esas evidencias se guardan en una carpeta institucional privada (cuenta \texttt{@uteq.edu.ec}) compartida con el docente, y en el repositorio se registra un manifiesto con nombre de archivo, fecha, duración y huella SHA-256 \citep{lopdp_2021}.
	\item \textbf{Uso de IA como apoyo.} Si se usa una herramienta de IA solo como apoyo (ortografía, error de compilación), se declara en una sección del documento: herramienta, versión, propósito y cómo se verificó lo obtenido.
\end{enumerate}

\seccion{Condiciones de entrega válida}

\begin{cajapiso}[title={Condiciones de entrega válida: su incumplimiento equivale a no entregar}]
	\begin{enumerate}[label=\textbf{P\arabic*.},leftmargin=24pt,itemsep=3pt]
		\item \textbf{PDF de una sola página con datos informativos y URL.} En el SGA se sube solamente un PDF de UNA SOLA página con los datos informativos (estudiante o estudiantes, tarea, fecha límite de entrega) y, en una sola línea, la URL del repositorio. Si el PDF tiene más de una página, si se suben otros archivos, si la URL no cabe en una sola línea o está rota, o si el repositorio no es accesible, la actividad formativa se considera \textbf{NO ENTREGADA} y el estudiante no puede presentar la sumativa vinculada, en la que obtiene CERO.
		\item \textbf{PDF reproducible.} Si el PDF no se genera correctamente a partir del \LaTeX{} al clonar el repositorio y compilar el archivo \texttt{.tex} siguiendo el \texttt{README.md} (compilador, orden de comandos, archivo principal y dependencias), la actividad formativa se considera \textbf{NO ENTREGADA} y el estudiante no puede presentar la sumativa vinculada, en la que obtiene CERO.
		\item \textbf{Plazo.} Si el contenido evaluable se sube después de la fecha y hora de cierre, la actividad formativa se considera \textbf{NO ENTREGADA} y el estudiante no puede presentar la sumativa vinculada, en la que obtiene CERO.
		\item \textbf{Integridad.} Plagio verificado, cualquiera de los indicios graves G1 a G6 de uso de IA, o contenido generado por IA presentado como propio: la actividad formativa se considera \textbf{NO ENTREGADA} y el estudiante no puede presentar la sumativa vinculada, en la que obtiene CERO, sin perjuicio de las acciones disciplinarias que correspondan.
	\end{enumerate}
\end{cajapiso}

\seccion{Integridad académica e indicios de uso de IA}

Todo trabajo debe ser de autoría propia. Se aplica el Código de Ética de la ACM \citep{acm_code_2018} y la sección 6.1 del encuadre.

\begin{cajapiso}[title={Indicios graves de uso de IA: cada uno, por sí solo, implica CERO (P4)}]
	\begin{enumerate}[label=\textbf{G\arabic*.},leftmargin=24pt,itemsep=3pt]
		\item \textbf{Referencias con datos que no corresponden a la obra real:} DOI que lleva a una obra con autores, título, año, revista, volumen o páginas diferentes; DOI, ISBN o URL inexistentes; obras que no existen.
		\item \textbf{Citas que la fuente no respalda:} afirmaciones, cifras o citas textuales atribuidas a una fuente que no las contiene, o cláusulas, secciones o tablas de una norma (IEEE 29148, ISO/\allowbreak{}IEC 25010) que no existen o no dicen lo afirmado.
		\item \textbf{Restos de un asistente de IA} en el documento, el código, los comentarios o los commits (por ejemplo, «Claro, aquí tienes\ldots», «Como modelo de lenguaje\ldots»).
		\item \textbf{Evidencias de campo fabricadas:} entrevistas, encuestas, participantes, transcripciones, actas o fotografías que no existen o no tienen respaldo original.
		\item \textbf{Resultados o artefactos inexistentes:} el texto describe modelos, matrices, pruebas o métricas que no están en el repositorio o no coinciden con él.
		\item \textbf{Incapacidad de explicar lo entregado:} el estudiante no puede justificar ni modificar, cuando se le pide, lo que escribió, modeló o especificó.
	\end{enumerate}
\end{cajapiso}

\begin{cajaaviso}[title={Indicios complementarios (se acumulan)}]
	\begin{enumerate}[label=\textbf{I\arabic*.},leftmargin=24pt,itemsep=3pt]
		\item Redacción uniforme de principio a fin, o estilo distinto al del estudiante en otras entregas o en su exposición oral.
		\item Uso excesivo de los mismos conectores y fórmulas («además», «cabe destacar», «es importante señalar», «en conclusión»).
		\item Contenido genérico: requisitos, actores o reglas aplicables a cualquier sistema, sin relación con lo elicitado.
		\item Entrega en uno o pocos commits masivos, o mensajes de commit idénticos o automáticos.
		\item Residuos de formato ajeno a \LaTeX{} (marcas de Markdown, emojis) o metadatos del PDF de otra herramienta.
		\item Estructura, redacción o errores muy similares entre estudiantes o equipos distintos.
		\item Traducción literal o mezcla de idiomas sin justificación.
	\end{enumerate}
	Con dos indicios complementarios se resta 1,5 puntos (nivel B) y con tres o más se restan 3,0 (nivel C). La declaración de uso de IA no exime de responsabilidad ni anula estos indicios.
\end{cajaaviso}

\seccion{Lista de cotejo para la retroalimentación}

La actividad no tiene calificación. El docente marca cada indicador y escribe una observación que el estudiante debe atender en la actividad sumativa vinculada. \textbf{L} = Logrado (cumple con la calidad descrita); \textbf{EP} = En proceso (presente pero incompleto, impreciso o con errores que se señalan); \textbf{NL} = No logrado (ausente, incorrecto o no verificable). Que un elemento esté presente no basta: se valora su calidad.

{\small
\begin{longtable}{|C{0.04\linewidth}|L{0.47\linewidth}|C{0.035\linewidth}|C{0.035\linewidth}|C{0.035\linewidth}|L{0.17\linewidth}|}
	\caption{Lista de cotejo de GA-SIM-01.}\\
	\hline
	\thdr{N.\textsuperscript{o}} & \thdr{Indicador y calidad esperada (Logrado)} & \thdr{L} & \thdr{EP} & \thdr{NL} & \thdr{Observación}\\
	\hline
	\endfirsthead
	\hline
	\thdr{N.\textsuperscript{o}} & \thdr{Indicador y calidad esperada (Logrado)} & \thdr{L} & \thdr{EP} & \thdr{NL} & \thdr{Observación}\\
	\hline
	\endhead
	1 & Priorizaciones individuales completas (MoSCoW y 100 puntos) por rol. & $\square$ & $\square$ & $\square$ & \\
	\hline
	2 & Lista acordada con justificación de los cambios respecto de las individuales. & $\square$ & $\square$ & $\square$ & \\
	\hline
	3 & Conflictos clasificados correctamente por tipo. & $\square$ & $\square$ & $\square$ & \\
	\hline
	4 & Estrategia de resolución identificada y evaluada. & $\square$ & $\square$ & $\square$ & \\
	\hline
	5 & Reflexión individual de cada integrante, propia y específica. & $\square$ & $\square$ & $\square$ & \\
	\hline
\end{longtable}
}

\begin{cajaclave}[title={Cómo se usa esta retroalimentación}]
	El docente devuelve la lista de cotejo en el SGA antes de la sesión siguiente. Cada indicador marcado como EP o NL debe corregirse antes de presentar GA-PE-02 y EV-PRAC-02; el docente verificará esa corrección al calificarla. Entregar esta formativa de forma válida es requisito para presentar GA-PE-02 y EV-PRAC-02.
\end{cajaclave}

% ============================================================
% ANEXOS: material que entrega el docente
% ============================================================
\clearpage
\seccion{Anexo A. Caso y requisitos para la simulación}

\subseccion{Caso: aplicación CacaoLote}

La Asociación Cacaotera Brisas del Quevedo (organización ficticia creada para esta actividad) agrupa a 85 pequeños productores de cacao fino de aroma de las parroquias rurales de Quevedo, provincia de Los Ríos. Cada semana, los socios llevan su grano seco al centro de acopio de la asociación, donde un bodeguero lo pesa, mide la humedad y anota la entrega en un cuaderno. Cuando se reúnen unos 2\,000 kg de grano de calidad similar, la directiva arma un lote y lo ofrece por teléfono a compradores intermediarios y exportadores. Este proceso genera reclamos: los socios no saben con certeza cuánto entregaron ni cuánto les corresponde cobrar, los cuadernos se pierden o se mojan, y los compradores piden cada vez más información sobre el origen de cada lote.

La directiva ha decidido contratar el desarrollo de una aplicación, llamada provisionalmente CacaoLote, para registrar las entregas en el centro de acopio, agrupar el grano en lotes, venderlos a compradores registrados y liquidar a cada socio. El centro de acopio tiene señal de internet intermitente y la mayoría de los socios usa teléfonos Android de gama baja.

Las condiciones son estrictas: la primera versión debe estar en uso en 3 meses, antes del inicio de la cosecha principal, y el presupuesto aprobado por la asamblea alcanza solo para unos 42 puntos de esfuerzo, es decir, alrededor del 60\,\% de los 69 puntos que suman los requisitos solicitados (Tabla~\ref{tab:simreq}). Por eso, su grupo deberá decidir qué entra en la primera versión, qué se pospone y qué se descarta. Los nombres, cifras y situaciones de este caso son ficticios.

\subseccion{Requisitos solicitados}

El esfuerzo se expresa en puntos relativos estimados por el equipo de desarrollo: un requisito de 8 puntos cuesta aproximadamente el cuádruple que uno de 2 puntos. La suma de los requisitos que su grupo incluya en la primera versión no debe superar 42 puntos.

{\small
\begin{longtable}{L{0.08\linewidth}L{0.56\linewidth}L{0.15\linewidth}C{0.10\linewidth}}
	\caption{Requisitos solicitados para la primera versión de CacaoLote}\label{tab:simreq}\\
	\thdr{Código} & \thdr{Enunciado} & \thdr{Tipo} & \thdr{Esfuerzo}\\
	\endfirsthead
	\thdr{Código} & \thdr{Enunciado} & \thdr{Tipo} & \thdr{Esfuerzo}\\
	\endhead
	R01 & El sistema permitirá registrar a cada socio y su finca (nombre, sector, superficie sembrada). & Funcional & 2\\
	\rowcolor{rowgray}R02 & El sistema permitirá registrar cada entrega en el centro de acopio con peso, humedad, calidad del grano y socio. & Funcional & 5\\
	R03 & El sistema permitirá agrupar entregas en lotes de venta, conservando qué socio aportó cada cantidad. & Funcional & 5\\
	\rowcolor{rowgray}R04 & El sistema permitirá publicar los lotes disponibles para que los compradores registrados presenten ofertas. & Funcional & 5\\
	R05 & El sistema calculará la liquidación de cada socio según el peso, la calidad y el precio del lote vendido. & Funcional & 8\\
	\rowcolor{rowgray}R06 & El sistema emitirá al socio un comprobante de entrega en PDF en el momento de la entrega. & Funcional & 3\\
	R07 & El sistema enviará automáticamente certificados de trazabilidad de cada lote al sistema informático del comprador exportador. & Funcional & 13\\
	\rowcolor{rowgray}R08 & Un socio sin experiencia previa con la aplicación podrá consultar sus entregas y su liquidación sin ayuda en menos de 3 minutos. & Calidad (capacidad de interacción) & 5\\
	R09 & El registro de entregas funcionará sin conexión a internet y se sincronizará automáticamente al recuperar la señal, sin pérdida de datos. & Calidad (fiabilidad) & 8\\
	\rowcolor{rowgray}R10 & El sistema almacenará la cédula, la ubicación de la finca y la cuenta bancaria de cada socio y las compartirá con los compradores del lote. & Funcional & 5\\
	R11 & La aplicación móvil deberá funcionar en teléfonos Android 10 o superior con 2 GB de memoria RAM. & Restricción & 2\\
	\rowcolor{rowgray}R12 & Solo el tesorero podrá aprobar liquidaciones, y toda modificación de entregas o liquidaciones quedará registrada con usuario, fecha y hora. & Calidad (seguridad) & 8\\
	\rowcolor{lightgreen}\multicolumn{3}{r}{\textbf{Total solicitado / presupuesto disponible}} & \textbf{69 / 42}\\
\end{longtable}
}

\subseccion{Plantilla 1. Priorización individual}

Complete esta plantilla de forma individual, desde su rol y antes de la negociación. Asigne a cada requisito una categoría MoSCoW (M: debe tener; S: debería tener; C: podría tener; W: no tendrá en esta versión) y reparta exactamente 100 puntos entre los 12 requisitos según la importancia que tienen para su rol (votación ponderada). Un requisito puede recibir 0 puntos.

Nombre: \rule{0.45\linewidth}{0.4pt} \quad Rol: \rule{0.3\linewidth}{0.4pt}

{\small
\begin{longtable}{C{0.12\linewidth}C{0.12\linewidth}C{0.12\linewidth}L{0.54\linewidth}}
	\caption{Registro de la priorización individual}\label{tab:simprio}\\
	\thdr{Requisito} & \thdr{MoSCoW} & \thdr{Puntos /100} & \thdr{Justificación breve}\\
	\endfirsthead
	\thdr{Requisito} & \thdr{MoSCoW} & \thdr{Puntos /100} & \thdr{Justificación breve}\\
	\endhead
	R01\rule{0pt}{7mm} & & & \\
	\rowcolor{rowgray}R02\rule{0pt}{7mm} & & & \\
	R03\rule{0pt}{7mm} & & & \\
	\rowcolor{rowgray}R04\rule{0pt}{7mm} & & & \\
	R05\rule{0pt}{7mm} & & & \\
	\rowcolor{rowgray}R06\rule{0pt}{7mm} & & & \\
	R07\rule{0pt}{7mm} & & & \\
	\rowcolor{rowgray}R08\rule{0pt}{7mm} & & & \\
	R09\rule{0pt}{7mm} & & & \\
	\rowcolor{rowgray}R10\rule{0pt}{7mm} & & & \\
	R11\rule{0pt}{7mm} & & & \\
	\rowcolor{rowgray}R12\rule{0pt}{7mm} & & & \\
	\rowcolor{lightgreen}\textbf{Total}\rule{0pt}{7mm} & & \textbf{100} & \\
\end{longtable}
}

\needspace{12\baselineskip}
\subseccion{Plantilla 2. Registro de conflictos de la negociación}

La persona con el rol de analista completa esta plantilla durante los 20 minutos de negociación. Anote cada desacuerdo relevante, el tipo de conflicto, los roles implicados, la estrategia con la que se resolvió y el resultado (por ejemplo, la categoría MoSCoW final o la reformulación acordada del requisito).

\begin{cajaclave}[title={Referencia rápida para completar la plantilla}]
	\textbf{Tipos de conflicto:} de datos (información distinta o incompleta), de intereses (objetivos diferentes), de valores (principios o criterios distintos), de relación (tensiones personales o de comunicación) y estructurales (presupuesto, plazo, autoridad o reglas fijas).

	\textbf{Estrategias de resolución:} acuerdo, compromiso (cada parte cede algo), votación, variante de solución (una alternativa nueva que satisface a las partes) y decisión por autoridad.
\end{cajaclave}

{\small
\begin{longtable}{L{0.24\linewidth}L{0.13\linewidth}L{0.14\linewidth}L{0.15\linewidth}L{0.20\linewidth}}
	\caption{Registro de conflictos y de su resolución}\label{tab:simconf}\\
	\thdr{Conflicto} & \thdr{Tipo} & \thdr{Roles implicados} & \thdr{Estrategia} & \thdr{Resultado}\\
	\endfirsthead
	\thdr{Conflicto} & \thdr{Tipo} & \thdr{Roles implicados} & \thdr{Estrategia} & \thdr{Resultado}\\
	\endhead
	\rule{0pt}{14mm} & & & & \\
	\rowcolor{rowgray}\rule{0pt}{14mm} & & & & \\
	\rule{0pt}{14mm} & & & & \\
	\rowcolor{rowgray}\rule{0pt}{14mm} & & & & \\
	\rule{0pt}{14mm} & & & & \\
	\rowcolor{rowgray}\rule{0pt}{14mm} & & & & \\
\end{longtable}
}

Lista final acordada (códigos y categoría MoSCoW): \rule{0.5\linewidth}{0.4pt}

Esfuerzo total de la lista final: \rule{0.12\linewidth}{0.4pt} puntos (máximo 42).

\bibliographystyle{unsrtnat}
\bibliography{referencias}

\end{document}
